% ID: SEL-D0005-OB-E03
% ESTADO: cotejada. Segunda pasada visual completa y borrador compilado revisados.
% Fuente: Selectividad/sel_2019_matematicas_aplicadas/Reserva A/Exámenes Mat Aplicadas CCSS - A y B.pdf
% SHA-256: d28871f7f203cf1be9e479e9d642cdb6b1623704a8b76a2b57cfa2d12950f293
% Página: 2; opción B; ejercicio 3.

Una determinada enfermedad puede estar provocada por una sola de las causas,
A, B o C. En el 35\,\% de los casos está provocada por A, en el 40\,\% por B
y en el 25\,\% por C.

Se sabe que el tratamiento de esta enfermedad requiere hospitalización en el
15\,\% de los casos si está provocada por A, en el 45\,\% si está provocada
por B y en un 20\,\% si está provocada por C. Se elige al azar una persona
afectada por esa enfermedad.
\begin{enumerate}[label=\alph*)]
  \item \textbf{(1.5 puntos)} \textquestiondown{}Cuál es la probabilidad de
  que necesite hospitalización?
  \item \textbf{(1 punto)} Si no necesita hospitalización,
  \textquestiondown{}cuál es la probabilidad de que la causa de la enfermedad
  sea C?
\end{enumerate}
